\documentclass[11pt]{article}
\usepackage[margin=1in]{geometry}
\usepackage[utf8]{inputenc}
\usepackage[T1]{fontenc}
\usepackage{lmodern}
\usepackage{amsmath,amssymb,amsthm}
\usepackage{booktabs}
\usepackage[title]{appendix}
\usepackage[colorlinks=true,linkcolor=blue,citecolor=blue,urlcolor=blue,pdfauthor={P. Graham},pdftitle={Competitive equilibrium can fail near equal budgets for three additive agents}]{hyperref}
\newtheorem{theorem}{Theorem}[section]
\newtheorem{corollary}[theorem]{Corollary}
\newtheorem{lemma}[theorem]{Lemma}
\theoremstyle{definition}
\newtheorem{definition}[theorem]{Definition}
\newtheorem{example}[theorem]{Example}
\theoremstyle{remark}
\newtheorem{remark}[theorem]{Remark}
\setlength{\emergencystretch}{3em}
\title{Competitive equilibrium can fail near equal budgets for three additive agents}
\author{P. Graham\\ \small\url{https://orcid.org/0009-0007-8493-1505}}
\date{November 5, 2026}
\begin{document}
\maketitle
\begin{abstract}
In a Fisher market with indivisible goods, each agent has a budget of
artificial money and a competitive equilibrium is an allocation of all the
goods with prices at which every agent gets a most preferred bundle within
budget. With exactly equal budgets one good and two agents already have no
equilibrium. Babaioff, Nisan and Talgam-Cohen proved that two agents with
additive preferences over any number of goods have a competitive
equilibrium whenever their budgets are sufficiently close to equal but not
equal, and asked whether this extends to more than two agents (their Open
Problem~1). It does not. We give three agents with strict additive
valuations over five goods such that for every $0<\delta<1/17$ no
competitive equilibrium exists at the budgets $(1-3\delta,1+2\delta,1+\delta)$,
nor at any budget vector near them. So every neighborhood of equal budgets
contains a set of positive measure on which no equilibrium exists. Adding
goods of small value, as in Segal-Halevi's Lemma~5, extends this to every
number of goods $m\ge5$. This answers the case of three additive agents and
$m\ge5$ goods that Segal-Halevi left open: the budgets with no competitive
equilibrium can have positive measure. The proof gives, for each allocation
of the goods, an exact Farkas (linear-programming duality) certificate that
no prices support it; an independent exact linear-programming computation
confirms it.

This work has not yet been reviewed by experts.
The proofs, the computer checks and the text were produced by AI models
(from OpenAI and Anthropic), working under the author's direction.
\end{abstract}
\medskip
\noindent\textit{Keywords:} competitive equilibrium, indivisible goods,
Fisher market, generic budgets, unequal entitlements, fair division.

\noindent\textit{MSC 2020:} 91B50, 91B32, 90C05.

\section{Introduction}

\subsection{Main result}
A set $M$ of $m$ indivisible goods is to be divided among $n$ agents. Agent
$i$ has an additive valuation $V_i$: a positive value $V_i(g)$ for each good
$g$, and $V_i(S)=\sum_{g\in S}V_i(g)$ for each bundle $S\subseteq M$. Agent
$i$ strictly prefers $S$ to $T$ when $V_i(S)>V_i(T)$, and the preferences are
\emph{strict} when the $2^m$ bundle values of each agent are pairwise
distinct. Each agent $i$ has a budget $t_i>0$. An \emph{allocation}
$X=(X_1,\dots,X_n)$ is a partition of $M$ into $n$ bundles, some possibly
empty; prices $p_g>0$ extend to bundles by $p(S)=\sum_{g\in S}p_g$.

\begin{definition}\label{def:ce}
A \emph{competitive equilibrium} (CE) at budgets $t=(t_1,\dots,t_n)$ is a
pair $(X,p)$ of an allocation and prices such that for every agent $i$:
(1)~$p(X_i)\le t_i$, and (2)~$p(Y)>t_i$ for every bundle $Y$ with
$V_i(Y)>V_i(X_i)$.
\end{definition}

This is Segal-Halevi's Definition~1 \cite{SH}, specialized to additive
preferences; Babaioff, Nisan and Talgam-Cohen \cite[Definition~2.1]{BNT}
use the same conditions with nonnegative prices. The theorem below holds
under either convention: the proof rules out even nonnegative prices with
weak inequalities in (2) (Lemma~\ref{lem:farkas}). Let $e=(1,1,1)$ denote
equal budgets.

\begin{table}[ht]
\centering
\begin{tabular}{lccccc}
\toprule
 & $v$ & $w$ & $x$ & $y$ & $z$\\
\midrule
agent 1 & 6 & 12 & 11 & 9 & 13\\
agent 2 & 14 & 4 & 1 & 8 & 20\\
agent 3 & 8 & 6 & 9 & 16 & 12\\
\bottomrule
\end{tabular}
\caption{The values $V_i(g)$ of the three agents for the five goods.}
\label{tab:values}
\end{table}

\begin{theorem}\label{thm:main}
Consider three agents with the additive valuations of Table~\ref{tab:values}
over the goods $M=\{v,w,x,y,z\}$; their preferences are strict. There is an
open set $U\subset\mathbb{R}^3_{>0}$, closed under multiplication by positive
scalars, such that no CE exists at any budget vector in $U$, and $U$
contains
\[
  e+\delta(-3,2,1)=(1-3\delta,\;1+2\delta,\;1+\delta)\qquad\text{for every }0<\delta<1/17 .
\]
Consequently every neighborhood of $e$, and of each positive multiple of
$e$, contains a set of budget vectors of positive Lebesgue measure at which
no CE exists.
\end{theorem}

For example, no CE exists at the budgets $(97,102,101)$ or $(17,22,21)$,
the multiples of $e+\delta(-3,2,1)$ for $\delta=1/100$ and $\delta=1/20$.
The theorem does not say that every budget vector near $e$ lacks a CE:
near $e$, the set $U$ is exactly the set of budgets with $t_1<t_3<t_2$ and
$t_1+t_2<2t_3$ (Remark~\ref{rem:cone}).

\begin{corollary}\label{cor:m}
For every $m\ge5$ there are three agents with strict additive valuations
over $m$ goods and a nonempty open set of budget vectors, closed under
positive scaling and with $e$ in its closure, at which no CE exists.
\end{corollary}

\subsection{The two questions}
Competitive equilibrium from equal incomes is a classical rule for fair
and efficient division \cite{Fol,Var}. With indivisible goods it can fail
even for one good and two agents, since whoever doesn't get the good can
afford it. Babaioff, Nisan and Talgam-Cohen credit Budish \cite{Bud} with initiating
``the study of competitive equilibrium with almost-equal budgets'', with
``existence results for an approximate such equilibrium notion'', and
``take the idea of almost-equal budgets a step further'', to exact
equilibria with generic budgets \cite[Section~1]{BNT}.
Their main result is \cite[Theorem~1.1]{BNT}: ``Consider 2 agents with
additive preferences over multiple items. If the budgets are sufficiently
close to being equal, but not precisely so, then a CE exists'' (in full,
\cite[Theorem~7.1]{BNT}). They then pose \cite[Open Problem~1]{BNT}:
\begin{quote}
Consider $n>2$ agents with additive preferences over multiple items. Does
there exist a CE when budgets are sufficiently close to being equal?
\end{quote}
Open Problem~1 appears from the second arXiv version of their paper
(September 25, 2018). The first (March 23, 2017) had asked more broadly
\cite[Section~1.4]{BNT1}: ``Does there always exist a CE for any number of players with
additive preferences and generic budgets?'' Segal-Halevi's Theorem~5 \cite{SH}
answers that question negatively with four agents, and
Theorem~\ref{thm:main} does so with three. Babaioff, Nisan and
Talgam-Cohen note that Segal-Halevi's four-agent result ``does not rule
out that for more than 2 agents, equilibrium exists for generic budgets
that are almost equal'' \cite[Section~1]{BNT}, where generic existence
means existence ``for every vector of budgets except for a zero-measure
subset'' \cite[Section~2.2.2]{BNT}. They also report a computer search: on
instances from the Spliddit website with additive preferences and more
than two agents, ``after perturbing the budgets to make them almost
equal, a CE was found'' \cite[Section~9.2 and Appendix~B.2]{BNT}. The
example here is not one of those instances.

Open Problem~1 can be read as asking for a CE at every unequal budget
vector near $e$, or at almost every one (generic existence).
Theorem~\ref{thm:main} answers both readings negatively: for the three
agents of Table~\ref{tab:values} there is no $\varepsilon>0$ such that a CE
exists at almost all budget vectors within $\varepsilon$ of $e$.

Segal-Halevi's Definition~3 \cite{SH} reads: ``Given integers $n\ge2$ and
$m\ge1$, a CE exists for almost all incomes if the set of income-vectors
for which no CE exists has a zero measure'', where by Definition~2 a CE
exists for an income vector $t$ if one exists ``for all
preference-profiles'' of the class. For additive preferences Segal-Halevi
proved that three agents and four goods have a CE for almost all incomes
\cite[Theorem~4]{SH}, and that with four or more agents and four or more
goods there are additive preferences and a positive-measure set of
incomes without a CE \cite[Theorem~5]{SH}. For monotone, non-additive
preferences, two agents and five goods may have no CE on a
positive-measure set of incomes \cite[Theorem~3.5]{BNT2},
\cite[Theorem~6]{SH}. Section~10.1 of \cite{SH} states: ``for $n=2$ or
$n=3$ agents with additive preferences, when there are $m\ge5$ goods, it is
open whether CE exists for almost all incomes.'' Wikipedia's list of
unsolved problems in fair division records as an open conjecture that
``When there are three agents, even with additive valuations, CE for
almost all incomes might not exist'' \cite{Wik}. Corollary~\ref{cor:m}
settles the case $n=3$: the answer is no for every $m\ge5$, as that
conjecture expects (Corollary~\ref{cor:questions}(b)). With the results of
\cite{SH} for $m\le4$ \cite[table in Section~1, Theorem~4]{SH}, this
completes the case of three additive agents, and $m=5$ is the least number
of goods for which a CE can fail on a positive-measure set of incomes.

What this paper adds is the example and its proof. The definitions, the
questions and the lifting argument for more goods are those of Babaioff,
Nisan and Talgam-Cohen and of Segal-Halevi. The difficult step was finding
an additive example: additivity rules out the preferences used for
monotone valuations (Segal-Halevi explains why the profile of
\cite[Theorem~6]{SH} is not additive \cite[Section~10.1]{SH}), and the
example here was found by computer search. This paper does not address two
additive agents with $m\ge5$ goods, which Segal-Halevi lists as open, or
Open Problem~1 for a fixed number of agents $n\ge4$.

\subsection{Proof overview}
At a fixed budget vector, deciding whether a CE exists is finite: there
are $3^5=243$ allocations, and for each one the CE conditions are linear
inequalities in the prices. We drop strictness, allowing $p\ge0$ and
$p(Y)\ge t_i$, and give for each allocation nonnegative multipliers whose
combination of these inequalities gives a linear inequality on the budgets
alone, $\sum_i w_it_i\ge0$ (Lemma~\ref{lem:farkas}). Where the budgets
violate it, that allocation has no CE prices. Along the ray
$e+\delta(-3,2,1)$ the inequality $\sum_iw_it_i<0$ is linear in $\delta$,
and the $243$ certificates together cover exactly $0<\delta<1/17$. Each
certificate excludes an open half-space of budgets, which gives the open
set $U$. Two kinds of obstruction appear. In the first, an agent prefers
part of another agent's bundle to its own, which forces a comparison of
budgets that fails on the ray (Example~\ref{ex:envy}). In the second, an
agent's preferred bundles together use the goods of other agents' bundles,
so their prices, each above that agent's budget, are bounded by the other
budgets (Example~\ref{ex:bottleneck}); an obstruction of this kind bounds
$\delta$ by $1/17$.

\section{Certificates}

\begin{lemma}\label{lem:farkas}
Let $X$ be an allocation and $t$ a budget vector. Let $P$ be the set of
pairs $(i,Y)$ with $V_i(Y)>V_i(X_i)$. Suppose $\alpha_i\ge0$ for each agent
and $\lambda_{i,Y}\ge0$ for $(i,Y)\in P$ satisfy, for every good $g$,
\begin{equation}\label{eq:cover}
  \sum_{i:\,g\in X_i}\alpha_i\;\ge\;\sum_{(i,Y)\in P:\,g\in Y}\lambda_{i,Y},
\end{equation}
and let $w_i=\alpha_i-\sum_{Y:(i,Y)\in P}\lambda_{i,Y}$. If
$\sum_iw_it_i<0$, then no prices $p\ge0$ satisfy $p(X_i)\le t_i$ for all
$i$ and $p(Y)\ge t_i$ for all $(i,Y)\in P$. In particular there is no CE
$(X,p)$ at $t$, nor at any budget vector $t'$ with $\sum_iw_it'_i<0$.
\end{lemma}

\begin{proof}
Suppose such $p$ exist. Multiply $t_i-p(X_i)\ge0$ by $\alpha_i$ and
$p(Y)-t_i\ge0$ by $\lambda_{i,Y}$, and add:
\[
  \sum_iw_it_i\;\ge\;\sum_g p_g\Bigl(\sum_{i:\,g\in X_i}\alpha_i-\sum_{(i,Y)\in P:\,g\in Y}\lambda_{i,Y}\Bigr)\;\ge\;0,
\]
by \eqref{eq:cover} and $p\ge0$, contradicting $\sum_iw_it_i<0$. A CE
$(X,p)$ at $t$ would satisfy these relaxed inequalities, since $p>0$ and
$p(Y)>t_i$ imply $p\ge0$ and $p(Y)\ge t_i$; so none exists, and the same
holds at any $t'$ with $\sum_iw_it'_i<0$.
\end{proof}

By Farkas' lemma \cite{Sch}, such multipliers exist whenever the relaxed
system has no solution, so nothing is lost by asking for them. We call
$w=(w_1,w_2,w_3)$ the \emph{weight} of the certificate. Two examples show
how the certificates read.

\begin{example}\label{ex:envy}
Let $X=(vwx,\,z,\,y)$, writing bundles as words. Agent~3 values $y$ at $16$
and $vx$ at $17$, so it prefers part of agent~1's bundle to its own. In a
CE, $p(vx)\le p(vwx)\le t_1$ while $p(vx)>t_3$, so $t_1>t_3$. The
certificate is $\alpha_1=\lambda_{3,vx}=1/2$, with weight
$(1/2,0,-1/2)$. Along the ray, $t_1-t_3=-4\delta<0$, so this allocation is
excluded for every $\delta>0$. Seventeen allocations are excluded by a
linear inequality between budgets that holds with equality at $e$
(Appendix~\ref{app:typeB}).
\end{example}

\begin{example}\label{ex:bottleneck}
Let $X=(w,\,vz,\,xy)$. Agent~1 values $w$ at $12$ and prefers $vx$ ($17$),
$vy$ ($15$) and $z$ ($13$). In a CE each of these costs more than $t_1$,
and since $vx+vy+2z$ uses each of $v,x,y,z$ as often as $2\,vz+xy$ does,
\[
  4t_1< p(vx)+p(vy)+2p(z)=2p(vz)+p(xy)\le 2t_2+t_3 .
\]
On the ray, $4t_1-2t_2-t_3=1-17\delta$, which is positive for
$\delta<1/17$. This certificate ($\alpha_2=2/7$, $\alpha_3=1/7$,
$\lambda_{1,vx}=\lambda_{1,vy}=1/7$, $\lambda_{1,z}=2/7$) is the one that
limits the theorem to $\delta<1/17$.
\end{example}

\section{Proof of Theorem~\ref{thm:main}}

Write $d=(-3,2,1)$. For each of the $3^5=243$ allocations $X$, a
certificate as in Lemma~\ref{lem:farkas}, with weight $w^X$, is listed in
the certificate file that accompanies this paper
(Section~\ref{sec:verification} says where); the counts and the bound
$1/17$ below are checked by the programs described there. Any direction
$d$ with $d_1<d_3<d_2$ and $d_1+d_2<2d_3$ would serve for small $\delta$
(Remark~\ref{rem:cone}); the certificates were computed for
$d=(-3,2,1)$. Each weight is of one of two kinds:
\begin{itemize}
\item[(A)] $w^X\!\cdot e<0$ ($226$ allocations). Then
$w^X\!\cdot(e+\delta d)=w^X\!\cdot e+\delta\,w^X\!\cdot d<0$ for
$0\le\delta<r_X$, where $r_X=-w^X\!\cdot e/w^X\!\cdot d$ if
$w^X\!\cdot d>0$ and $r_X=\infty$ otherwise.
\item[(B)] $w^X\!\cdot e=0$ and $w^X\!\cdot d<0$ ($17$ allocations). Then
$w^X\!\cdot(e+\delta d)<0$ for every $\delta>0$.
\end{itemize}
The least $r_X$ is $1/17$, attained by Example~\ref{ex:bottleneck}.

\begin{proof}[Proof of Theorem~\ref{thm:main}]
The preferences are strict and the valuations strictly monotone
(Section~\ref{sec:verification}). Let
\[
  U=\{t\in\mathbb{R}^3_{>0}:\ w^X\!\cdot t<0\ \text{for all }243\text{ allocations }X\}.
\]
It is a finite intersection of open half-spaces with $\mathbb{R}^3_{>0}$,
so it is open, and it is closed under positive scaling. At $t\in U$, every
allocation $X$ is excluded by Lemma~\ref{lem:farkas} with its certificate,
so no CE exists. By (A) and (B), $e+\delta d\in U$ for every
$0<\delta<1/17$. Given a positive multiple $ce$ and $\varepsilon>0$, choose
$\delta$ with $0<\delta<1/17$ and $c\delta\lVert d\rVert<\varepsilon$; then $c(e+\delta d)\in U$
is within $\varepsilon$ of $ce$, and since $U$ is open, a small ball around
it lies in $U$ and in the $\varepsilon$-ball around $ce$. That ball has
positive measure.
\end{proof}

\begin{remark}\label{rem:cone}
Near $e$, the set $U$ is cut out by the kind-(B) inequalities alone. The
$243$ certificates have $51$ distinct weights up to positive scaling, four
of kind (B). These four give the inequalities $t_1<t_3$, $t_3<t_2$,
$t_1<t_2$ and $t_1+t_2<2t_3$ (Appendix~\ref{app:typeB}), and the kind-(A)
inequalities hold strictly at $e$. So there is a neighborhood $N$ of $e$
in which $U\cap N$ is the intersection of $N$ with the open polyhedral
cone defined by $t_1<t_3<t_2$ and $t_1+t_2<2t_3$ (the inequality
$t_1<t_2$ follows from these).
\end{remark}

\begin{remark}
The certificates use only which bundles each agent strictly prefers to
which. So the theorem holds for every additive profile that ranks the
bundles as Table~\ref{tab:values} does, an open set of valuation profiles.
\end{remark}

\section{More goods, and the two questions}

\begin{proof}[Proof of Corollary~\ref{cor:m}]
We follow Segal-Halevi's proof of \cite[Lemma~5(a)]{SH}, which adds a
``low-value item'' that changes no comparison between the original
bundles. Let $k=m-5$ and give each agent value $2^{-j}$ for the new good
$g_j$, $j=1,\dots,k$. Any set $Z$ of new goods has value less than $1$.
The original values are integers and the preferences are strict, so
$V_i(W)>V_i(Y)$ implies $V_i(W)\ge V_i(Y)+1>V_i(Y\cup Z)$ for original
bundles $W,Y$. Bundles of the larger market with different original parts
therefore have different values, and bundles with the same original part
differ by distinct sums of distinct powers of $2$; so the preferences stay
strict. Let $t\in U$ and suppose $(X,p)$ is a CE of the larger market at
$t$. Put $X'_i=X_i\cap M$. Then $p(X'_i)\le p(X_i)\le t_i$, and if
$W\subseteq M$ has $V_i(W)>V_i(X'_i)$ then $V_i(W)>V_i(X_i)$, so
$p(W)>t_i$. Thus $(X',p|_M)$ is a CE of the five-good market at $t$,
contradicting Theorem~\ref{thm:main}. So no CE exists at any $t\in U$.
\end{proof}

\begin{corollary}\label{cor:questions}
(a) The answer to \cite[Open Problem~1]{BNT} is no: for the three additive
agents of Table~\ref{tab:values} there is no $\varepsilon>0$ such that a CE
exists for almost all budget vectors within $\varepsilon$ of equal budgets.
(b) For three agents with additive preferences and $m\ge5$ goods, a CE
need not exist for almost all incomes in the sense of
\cite[Definition~3]{SH}.
\end{corollary}

\begin{proof}
Both follow from Theorem~\ref{thm:main} and Corollary~\ref{cor:m}: one
profile has no CE on a set of positive measure within every
$\varepsilon$-ball around $e$. Normalizing budgets to sum to $1$, as
\cite[Section~2.1]{BNT} does, changes nothing: the coordinates of $d$ sum
to $0$, and $U$ is closed under scaling.
\end{proof}

\section{Computational verification}\label{sec:verification}

Three short programs in exact rational arithmetic check the theorem from
the market and the certificate file. The programs, the certificate file
\texttt{data/certificate.json} and the command \texttt{sh checks/check.sh}
that runs them are in the release accompanying this paper,
\url{https://github.com/6/clank-lab/tree/near-equal-budgets-v1/results/near-equal-budgets};
its README describes the file format and the requirements. The first checks the hypotheses: the values are positive
integers, each agent's $32$ bundle values are distinct, and the market in
the certificate file is that of Table~\ref{tab:values}. The second reads
the certificate for each of the $243$ allocations; checks that every
$\lambda_{i,Y}$ is attached to a bundle $Y$ that agent $i$ strictly prefers
to $X_i$, that the multipliers are nonnegative and satisfy
\eqref{eq:cover}, and that the weight is of kind (A) or (B); and computes
the least $r_X$, which is $1/17$. It runs in well under a second.

The third program does not read the certificates. For each allocation it
solves, with its own exact two-phase simplex method, the linear program in
the prices $p\ge0$ and $\delta\ge0$ with $p(X_i)\le 1+\delta d_i$ and
$p(Y)\ge1+\delta d_i$ for every bundle $Y$ agent $i$ strictly prefers to
$X_i$. The feasible values of $\delta$ form a closed interval, and the
program computes its left end $\delta_X$. It finds $214$ allocations
infeasible for every $\delta\ge0$, $17$ feasible only at $\delta=0$ (the
allocations of kind (B)), and $12$ whose feasible interval has left end
$\delta_X\ge1/17$, the least $\delta_X$ being exactly $1/17$. So the two
programs agree on the bound $1/17$, and by Farkas' lemma no certificate as
in Lemma~\ref{lem:farkas} exists at $\delta=1/17$ for at least one
allocation. Whether a CE exists at some
$\delta\ge1/17$ is not decided here. The program takes about four
minutes on one core.

As sanity checks, the second program rejects the certificates with one
multiplier doubled. With agent~1's value for $v$ changed from $6$ to $1$,
the second program rejects the original certificates, and the third finds
an allocation feasible for all $\delta$ in $[0,2/11]$. The example was found by a randomized search with
floating-point linear programs and reduced to small integers; the search
is not part of the proof.

\begin{appendices}

\section{The seventeen allocations of kind (B)}\label{app:typeB}

Each allocation below, written $(X_1,X_2,X_3)$, is excluded at budgets $t$
by the single inequality shown. The middle column names the bundles $Y$
whose prices the certificate compares with the budgets; each is strictly
preferred by the named agent to its own bundle. For instance, in the first
row agent~3 prefers $vx$ ($V_3=17>16$), which lies inside $X_1$. All seventeen are excluded on the
ray for every $\delta>0$, since $t_1-t_3=-4\delta$, $t_3-t_2=-\delta$,
$t_1-t_2=-5\delta$ and $t_1+t_2-2t_3=-3\delta$.

\begin{center}
\begin{tabular}{lll}
\toprule
allocation & preferred bundles used & excluded where\\
\midrule
$(vwx,z,y)$, $(vx,wz,y)$ & agent 3: $vx$ & $t_1<t_3$\\
$(wy,z,vx)$ & agent 3: $wy$ & $t_1<t_3$\\
$(vw,xz,y)$ & agent 2: $wz$; agent 3: $vx$ & $t_1<t_3$\\
$(wxz,v,y)$, $(wz,vx,y)$, $(wz,v,xy)$, $(xz,vw,y)$ & agent 2: $z$ & $t_1<t_2$\\
$(z,vw,xy)$, $(xz,v,wy)$, $(z,vx,wy)$, $(wz,y,vx)$ & agent 2: $z$ & $t_1<t_2$\\
$(wx,v,yz)$ & agent 2: $z$; agent 3: $wxy$ & $t_1<t_2$\\
$(wx,z,vy)$ & agent 2: $vy$ & $t_3<t_2$\\
$(vw,z,xy)$, $(vx,z,wy)$ & agent 1: $wx$; agent 2: $vy$ & $t_3<t_2$\\
$(wx,vz,y)$ & agent 3: $vx$, $wz$ & $t_1+t_2<2t_3$\\
\bottomrule
\end{tabular}
\end{center}

For example, for $(vw,z,xy)$: agent~1 prefers $wx$ ($23>18$) and agent~2
prefers $vy$ ($22>20$); together $wx$ and $vy$ use the goods of
$X_1\cup X_3=vwxy$ once each, so in a CE $t_1+t_2<p(wx)+p(vy)=p(vw)+p(xy)\le
t_1+t_3$. For $(wx,vz,y)$, agent~3 prefers $vx$ ($17>16$) and $wz$
($18>16$), which together use the goods of $X_1\cup X_2$ once each, so a
CE needs $2t_3<p(vx)+p(wz)=p(wx)+p(vz)\le t_1+t_2$.

\end{appendices}

\begin{thebibliography}{9}

\bibitem{BNT}
M.~Babaioff, N.~Nisan and I.~Talgam-Cohen, \emph{Competitive equilibrium
with indivisible goods and generic budgets}, Math. Oper. Res. \textbf{46}
(2021), 382--403;
\href{https://doi.org/10.1287/moor.2020.1062}{doi:10.1287/moor.2020.1062};
\href{https://arxiv.org/abs/1703.08150v2}{arXiv:1703.08150v2}, the version
cited.

\bibitem{BNT1}
M.~Babaioff, N.~Nisan and I.~Talgam-Cohen, \emph{Competitive equilibria
with indivisible goods and generic budgets},
\href{https://arxiv.org/abs/1703.08150v1}{arXiv:1703.08150v1} (March 23,
2017), the first version of \cite{BNT}.

\bibitem{BNT2}
M.~Babaioff, N.~Nisan and I.~Talgam-Cohen, \emph{Competitive equilibrium
with generic budgets: beyond additive},
\href{https://arxiv.org/abs/1911.09992v1}{arXiv:1911.09992v1} (2019).

\bibitem{Bud}
E.~Budish, \emph{The combinatorial assignment problem: approximate
competitive equilibrium from equal incomes}, J. Political Econ.
\textbf{119} (2011), 1061--1103;
\href{https://doi.org/10.1086/664613}{doi:10.1086/664613}.

\bibitem{Fol}
D.~K.~Foley, \emph{Resource allocation and the public sector}, Yale
Econ. Essays \textbf{7} (1967), 45--98.

\bibitem{Sch}
A.~Schrijver, \emph{Theory of Linear and Integer Programming}, Wiley,
Chichester, 1986.

\bibitem{SH}
E.~Segal-Halevi, \emph{Competitive equilibrium for almost all incomes:
existence and fairness}, Auton. Agents Multi-Agent Syst. \textbf{34}
(2020), article 26;
\href{https://doi.org/10.1007/s10458-020-09444-z}{doi:10.1007/s10458-020-09444-z};
\href{https://arxiv.org/abs/1705.04212v6}{arXiv:1705.04212v6}, the version
cited.

\bibitem{Var}
H.~R.~Varian, \emph{Equity, envy, and efficiency}, J. Econom. Theory
\textbf{9} (1974), 63--91;
\href{https://doi.org/10.1016/0022-0531(74)90075-1}{doi:10.1016/0022-0531(74)90075-1}.

\bibitem{Wik}
Wikipedia, \emph{List of unsolved problems in fair division}, section
``Competitive equilibrium for almost all incomes'',
\url{https://en.wikipedia.org/wiki/List_of_unsolved_problems_in_fair_division},
revision 1360704274, accessed November 5, 2026.

\end{thebibliography}
\end{document}
